\documentclass[10pt,a4paper]{article}
\usepackage{amsmath,amssymb,amsthm,graphicx}
% Minimal retained preamble for lane-p2 cards taken from
% sources/ahl-16-fixed-points-2019/editorial-preamble.tex
% (doi:10.5802/ahl.16, Annales Henri Lebesgue 2 (2019) 149-185, CC BY 4.0).
%
% Copied verbatim from the editorial preamble: the package list, the
% theorem declarations under the shared counter and the macro block, so that
% every macro and environment of the retained spans is available.
%
% Omitted and why: the class line (the assembler supplies the class); the
% a4paper geometry line and the fontspec block with its explicit lmroman10
% files (encodings and fonts fight the forced xelatex engine, and the
% assembler supplies the page setup).
\usepackage{amsmath,amssymb,amsthm,mathtools,graphicx,xcolor,hyperref,xurl,xspace,ifthen}
\newtheorem{theo}{Theorem}[section]
\newtheorem{lemm}[theo]{Lemma}
\newtheorem{theorem}[theo]{Theorem}
\newtheorem*{remas*}{Remarks}
\providecommand{\og}{«\nobreakspace}
\providecommand{\fg}{\nobreakspace»}
\newtheorem{lem}[theo]{Lemma}
\newtheorem{lemma}[theo]{Lemma}
\newtheorem{prop}[theo]{Proposition}
\newtheorem{coro}[theo]{Corollary}
\newtheorem{corol}[theo]{Corollary}
\newtheorem{conj}[theo]{Conjecture}
\newtheorem{Proposition}[theo]{Proposition}
\newtheorem{theoreme}[theo]{Theorem}
\theoremstyle{definition}
\newtheorem{defi}[theo]{Definition}
\newtheorem{exam}[theo]{Example}
\newtheorem{exem}[theo]{Example}
\newtheorem{rema}[theo]{Remark}
\newtheorem{nota}[theo]{Notation}
\newtheorem*{rema*}{Remark}
\newtheorem*{defi*}{Definition}
\newtheorem*{theo*}{Theorem}
\newtheorem*{lemm*}{Lemma}
\newtheorem*{prop*}{Proposition}
\newcommand{\enoncename}{}
\newtheorem{corpusenonce}[theo]{\enoncename}
\newenvironment{enonce}[2][]{\renewcommand{\enoncename}{#2}\begin{corpusenonce}}{\end{corpusenonce}}
\newenvironment{enonce*}[2][]{\par\medskip\noindent\textbf{#2.}\itshape}{\par\medskip}
\newenvironment{proofc}[1][\proofname]{\begin{proof}[#1]}{\end{proof}}
\newcommand{\Mc}{\mathcal{M}}
\newcommand{\pointir}{.\ ---\ }
\providecommand{\mathof}[1]{\mathrm{#1}}

\numberwithin{equation}{section}
\numberwithin{figure}{section}
\numberwithin{table}{section}
\def\endappendix{}
%~Mouliné par MaN_auto v.0.21.4 2019-04-30 09:02:09


\usepackage[all]{xy}

\hyphenation{auto-morphism
every
Roelcke
peri-od-ic
peri-od-icy
rep-re-sen-ta-tion
}








\newcommand{\A}{\mathcal A}
\newcommand{\N}{\mathbb{N}}
\newcommand{\Z}{\mathbb{Z}}
\newcommand{\Q}{\mathbb{Q}}
\newcommand{\K}{\mathcal{K}}
\newcommand{\CC}{\mathbb{C}}

\newcommand*{\mA}{\mathbf{A}}
\newcommand*{\mB}{\mathbf{B}}
\newcommand*{\mC}{\mathbf{C}}
\newcommand*{\mD}{\mathbf{D}}
\newcommand*{\mE}{\mathbf{E}}
\newcommand*{\mF}{\mathbf{F}}
\newcommand*{\mH}{\mathbf{H}}
\newcommand*{\mZ}{\mathbf{Z}}
%\newcommand*{\mA}{\mathbf{A}}
%\newcommand*{\mA}{\mathbf{A}}

\newcommand{\restrict}[2]{#1 \upharpoonright #2}

\newcommand{\e}{\varepsilon}

\DeclareMathOperator{\Aut}{Aut}
\DeclareMathOperator{\Stab}{Stab}
\DeclareMathOperator{\Age}{Age}

\newcommand{\actson}{\curvearrowright}

\newcommand{\RUCB}{\mathrm{RUC}_{b}}
\newcommand{\ROB}{\mathrm{Ro}_{b}}
\newcommand{\WAP}{\mathrm{WAP}}
\newcommand{\AP}{\mathrm{AP}}

\graphicspath{{./figures/}}

\newcommand*{\mk}{\mkern -1mu}
\newcommand*{\Mk}{\mkern -2mu}
\newcommand*{\mK}{\mkern 1mu}
\newcommand*{\MK}{\mkern 2mu}

\hypersetup{urlcolor=purple, linkcolor=blue, citecolor=red}

\newcommand*{\romanenumi}{\renewcommand*{\theenumi}{\roman{enumi}}}
\newcommand*{\Romanenumi}{\renewcommand*{\theenumi}{\Roman{enumi}}}
\newcommand*{\alphenumi}{\renewcommand*{\theenumi}{\alph{enumi}}}
\newcommand*{\Alphenumi}{\renewcommand*{\theenumi}{\Alph{enumi}}}
\let\oldexists\exists
\renewcommand*{\exists}{\mathrel{\oldexists}}
\let\oldforall\forall
\renewcommand*{\forall}{\mathrel{\oldforall}}


\usepackage{cleveref}
\crefname{originalproposition}{original proposition}{original propositions}
\Crefname{originalproposition}{Original proposition}{Original propositions}
\makeatletter
\newlabel{prop:RoelckeRP}{{4.5}{source}{}{sourcepointer-1}{}}
\newlabel{prop:RoelckeRP@cref}{{[originalproposition][0][]4.5}{[1][0][]source}{}{}{}}
\makeatother
\title{A finite coloring is a Roelcke coloring exactly when it depends only on the pattern with one finite substructure}
\author{Lionel Nguyen Van Thé}
\date{doi:10.5802/ahl.16; CC BY 4.0}
\begin{document}
\maketitle
\noindent\textit{Editorial note.} The counted unit is the article's own Proposition prop:Roelcke coloring at source lines 515-517 together with its proof at source lines 523-525. That proof environment is shared with the statement that follows it, prop:RoelckeRP at lines 519-521; the block proves the counted statement first and then continues with the neighbour, and the neighbour is declared as a navigation pointer. This is the pairing the gate reported as clean, and the two statements are not interderivable: one describes a single coloring, the other an equivalence quantified over all finite substructures, embeddings and colorings of patterns. The source is the publisher version of the article in Annales Henri Lebesgue 2 (2019), pages 149-185, DOI 10.5802/ahl.16, released under CC BY 4.0.\par\medskip
\noindent\textit{Editorial note.} One earlier note of this lane called this statement uncountable on the ground that it has no proof of its own. That reading was wrong for this statement: the block at lines 523-525 does carry its argument, and the article prints the two statements as separate propositions because it uses each on its own. The recorded pairing makes the sharing explicit, and the neighbour is left as a pointer.\par\medskip
\noindent\textit{Editorial note.} The licence evidence used here is the observation file kept with the source, license-observation.json, which records the licence as CC BY 4.0, points at the inherited publisher-version evidence and states that the landing page displays no licence line; the underlying evidence files are the same ones the earlier cards of this source rely on. No acceptance record is written anywhere.\par\medskip
\noindent\textit{Editorial note.} The retained preamble is this source's editorial preamble with three omissions only: the class line, which the assembler supplies; the page geometry; and the fontspec block with its explicit font files, which would fight the forced xelatex engine. Everything else is retained verbatim; the card copies the shared slice kept at focused-seeds/logic/P2-ahl16-retained-preamble.tex.\par\medskip
\noindent\textit{Editorial note.} Neither the retained statement nor the retained proof contains a citation, so this card needs no bibliography entry.\par\medskip
\noindent\textit{Editorial note.} The census package marked this candidate free of missing citation keys and listed the statement span as 515-517 and the proof span as 523-525; both were re-read from the source before the card was written, together with the boundary of the shared block.\par\medskip
\noindent\textit{Editorial note.} Revision r002 of this card aligns the licence evidence pointer with the source observation file kept with the source, sources/ahl-16-fixed-points-2019/license-observation.json and the recorded version string is aligned verbatim with it. The counted statement, the proof spans and the navigation pointers are unchanged; only the pointer, and where it differed the recorded version string, moved.\par\medskip
\noindent\textit{Allowed standard prerequisites (editorial).}\begin{itemize}
\item Roelcke colorings and joint embedding patterns. For a Fraisse structure the article attaches to each finite substructure and each joint embedding of a finite substructure into the structure the pattern of a copy, and a finite coloring of the copies is a Roelcke coloring when it is uniformly continuous for the Roelcke uniformity of the automorphism group, which the article shows is the same as depending only on those patterns. The author's own definitions and the surrounding results are the ones used by the counted statement and proof.. The counted proposition is the combinatorial description of a single Roelcke coloring, and the article uses it to pass between a coloring of copies and a coloring of the patterns of one joint embedding.
\end{itemize}\medskip
\ifdefined\hypertarget\hypertarget{sourcepointer-1}{}\fi
\noindent\textit{Source reference.} Original proposition 4.5: sources/ahl-16-fixed-points-2019/body-01.tex, source lines 519--521. Navigation pointer for the neighbouring statement whose proof is the same block retained here; the block proves that statement after it has finished this one, and the retained span ends before that argument begins, so the reference is declared as a pointer carrying the number the article prints for that statement, that is the section together with the place of the statement inside it, read off an enumeration of every numbered theorem-like environment of the source in order: it is 4.5; the cited statement is an allowed input and none of it is reproduced inside the span.\par\medskip
\begin{prop}\label{prop:Roelcke coloring}
Let $\mF$ be a Fra\"{\i}ss\'e structure, $\mA\in \Age(\mF)$ and $\chi$ be a finite coloring of $\binom{\mF}{\mA}$. Then $\chi\in \ROB(\Aut(\mF))$ iff there is a finite substructure $\mZ$ of $\mF$ such that $\chi(a)$ depends only on $[a,z]$, where $z$ is the identity embedding.
\end{prop}

\begin{proof}
Assume that $\mF$ has the Ramsey property for colorings in $\ROB(\Aut(\mF))$, and fix $\mA, \mB, \mZ\in \Age(\mF)$, $z\in \binom{\mF}{\mZ}$, $\gamma$ a finite coloring of the joint embedding patterns of $\mA$ and $\mZ$. Then the coloring defined on $\binom{\mF}{\mA}$ by $a\mapsto \gamma([a,z])$ is in $\ROB(\Aut(\mF))$ by Proposition~\ref{prop:Roelcke coloring}. The conclusion follows. The converse is an immediate consequence of the following easy fact: if $\mathcal C = \{ \chi_{i}:i\leq n\} \subset \ROB(\Aut(\mF))$ is a finite set of finite colorings, which we may assume to be colorings of $\binom{\mF}{\mA}$, Proposition~\ref{prop:Roelcke coloring} guarantees that each $\chi_{i}$ is associated to some finite $\mZ^{i}\subset \Age(\mF)$ and finite coloring $\gamma_{i}$ of the joint embedding patterns of $\mA$ and $\mZ^{i}$. Then, we see that the hypothesis applied to $\mZ=\bigcup_{i\leq n}z^{i}(\mZ^{i})$ and $\gamma$ defined by $\gamma([a,z]) = (\gamma_{i}([a,z^{i}]))_{i\leq n}$ provides $b\in \binom{\mF}{\mB}$ so that for every $i\leq n$, the coloring $a\mapsto \gamma_{i}([a,z^{i}])$ is constant on $\binom{b(\mB)}{\mA}$.
\end{proof}

\end{document}
